\section{Reduction and angular levels}

Failure for every $C$ would give a sequence with $n/(d^4\log(d+1))\to\infty$. Bounded $d$ is excluded by the fixed-cap theorem, as finitely many caps are then involved. The full cap $r=d\ge2$ has simple negative roots by the binomial multiplier theorem. Fix the absolute near-saturation threshold $h_0$ from the upper saddle estimate first. The interior estimate then has fixed positive deficit parameter $\eta=h_0$, and its cutoff constant can be chosen after $h_0$ is fixed. Pass to $d\to\infty$ and $r<d$. Put $T=nr/d$, $D=\lfloor T\rfloor$, $b=d^2/n$ and $\nu=(d-1)/2$. We have $b\to0$ and $D/d\to\infty$. Every integer quota between 0 and $D$ is feasible. For instance, split $dk$ slots among the $n$ rows in groups of at most $r$ and fill the slots successively with the repeating sequence of columns $1,\ldots,d$. \href{https://github.com/ProgramComputer/open-math-research/blob/f7e10f84dd21fec44ca58c60bae0879d12ebd89b/topics/balanced-arrays/finite-linear-reduction.md}{Full-cap result}

At row tilt $w>0$ put $H(w)=\sum_{s=0}^r\binom ds w^s$, with normalized size law $S$, mean $\mu$, variance $V$, and $h=r-\mu$. Write $w_x$ for the tilt with $y(w_x)=x$, where $y=n\mu/d$. This covers $0<x<T$ since $\partial_{\log w}\mu=V>0$. For $t_*=\pi/d$ and $G(s)=s(d-s)/(d(d-1))$, define

\begin{equation*}
\lambda_t=\frac{\mathbb E[G(S)e^{itS}]}{\mathbb E e^{itS}},
\qquad \Psi(w)=n\arg H(we^{it_*})-\nu\arg\lambda_{t_*}.
\end{equation*}

At $0\le t\le t_*$ both positive-weight characteristic sums lie in the closed sector from zero angle to $rt<\pi$, so the arguments continued from zero exist. The limiting values of $\Psi$ at zero and infinity are $-\nu t_*$ and $T\pi$. The lower limit uses the first nonzero degree of the $G$-weighted polynomial, and the upper limit uses $G(r)>0$. For $1\le k<D$ define $v_k$ as the first radius attaining $\Psi=k\pi$. Continuity gives strictly increasing $v_k$.

The following estimate controls the locations of the levels with an unbounded cap:

\begin{equation}
|\Psi(w)/\pi-y|\le C_1 d+\frac{C_1}{\sqrt d}\min(y,T-y).
\tag{M1}
\end{equation}

For a coordinate indicator $X_1$ and $p=\mu/d$, exchangeability gives $V=d\,\mathrm{Cov}(S,X_1)=dp(1-p)(\mathbb E[S\mid X_1=1]-\mathbb E[S\mid X_1=0])$. Given $X_1=1$ or $X_1=0$, the number of other successes has the same $(d-1)$-trial tilted binomial law truncated at $r-1$ or $r$, respectively. The former truncation has smaller mean, so the conditional difference for $S$ is at most one and $V\le\mu(1-\mu/d)\le d/4$. Also $V\le m(1+m)$ for $m=\mu$ and for $m=h$. To see this for any nonnegative integer variable $X$ with log-concave mass on $[0,r]$ and mean $m$, write $F_j=\mathbb P(X\ge j)$. Its ratios $F_{j+1}/F_j$ decrease since $F_j/\mathbb P(X=j)$ is the sum of products of consecutive decreasing mass ratios. Hence $F_{i+j}\le F_iF_j$ and

\begin{equation*}
\mathbb E[X(X-1)]/2=\sum_{i,j\ge1}F_{i+j}\le m^2 .
\end{equation*}

Apply this to $S$ and $r-S$. Combining the two variance bounds at $m=\sqrt d$ gives $V\le C\sqrt d\min(\mu,h)$. Now $|\mathbb E e^{it(S-\mu)}-1|\le V t^2/2$ by Taylor with centered linear term. This is $O(1/d)$ on $[0,t_*]$, whence $|\arg H(we^{it_*})-\mu t_*|\le C V/d^2$. Finally $|\arg\lambda_{t_*}|\le\pi$ by the sector property. Substituting gives (M1). \hyperref[sec:interior-saddle]{Uniform interior saddle and variance bound}

Consequently, at $w=v_k$, bounded $k/d$ implies $y=O(d)$, and bounded $(D-k)/d$ implies $T-y=O(d)$. If both indices divided by $d$ tend to infinity, so do $y/d$ and $(T-y)/d$, by applying (M1) to either hypothetical bounded subsequence.

\section{The determinant phase at the ends}

Suppose $k/d\to c$ with $0<c<\infty$. By (M1), $\mu=O(b)$ at the crossing. The size law then has $\mathbb E|S-\mu|^3=O(\mu)$, and the law reweighted by $G$ has expected distance $O(\mu)$ from size one. Indeed $dw=O(\mu)$, $p_s\le(dw)^s/s!$, $p_1\asymp\mu$, and $G(s)/G(1)\le s$. Thus $\arg\lambda_{t_*}=t_*+O(\mu t_*)$. The centered characteristic argument at $t_*$ is $O(\mu t_*^3)$ by the vanishing linear term and the third absolute moment. Since $n\mu t_*^3=O(y/d^2)=o(1)$, the level equation gives $y(v_k)=k+1/2+o(1)$.

Suppose instead $(D-k)/d\to c$ with $0<c<\infty$. We have $h=O(b)$ at the crossing. For $J=r-S$ let $p'_j=\mathbb P(J=j)$. By log-concavity all forward ratios are at most $p'_1/p'_0\le h/(1-h)$. It follows that $\mathbb E|S-\mu|^3=O(h)$, and that the $G$-weighted mean distance from $r$ is $O(h)$, using $G(r-j)/G(r)\le1+j$. We obtain $\arg\lambda_{t_*}=O(h t_*)$ and a centered characteristic argument $O(h t_*^3)$. Thus $y(v_k)=k+o(1)$.

At real $x$ within bounded distance of either such $k$, the variance is asymptotic to $dx/n$ at the lower end and to $d(T-x)/n$ at the upper end. At the lower end this follows from the small-mean factorial moments. At the upper end it follows from the geometric domination of $J$, which gives $\mathbb E J(J-1)=O(h^2)$. Differentiation yields $\frac{\mathrm d}{\mathrm dx}(d\log w_x)=b/V_x$, tending to $1/c$ uniformly on each bounded window. Integrating from $k$ to $y(v_k)$ gives the two expansions

\begin{equation}
d\log v_k=d\log w_k+d/(2k)+o(1)\quad(k/d\to c),
\qquad d\log v_k=d\log w_k+o(1)\quad((D-k)/d\to c).
\tag{M2}
\end{equation}

\section{Signs and degree}

Write $t_k=\sqrt{a_{k-1}/a_{k+1}}$ and $\Gamma_k=a_k^2/(a_{k-1}a_{k+1})$. In either finite positive overlap above, endpoint coefficient estimates show that $\log\Gamma_{k+j}\to1/c$ for every fixed integer $j$ and $\log t_k=d\log v_k+o(1)$. The lower endpoint proof compares with independent columns. The upper endpoint proof conditions on the complete row-size profile, including at exact saturation. At distances with $k/d\to0$ or $(D-k)/d\to0$, they show $\Gamma_k\to\infty$. \hyperref[sec:nearly-empty]{Nearly empty rows} \enspace \hyperref[sec:upper-quota]{Nearly saturated coefficients}

By the orientation inequality the $a_k$ are strictly log-concave at internal quotas. If $\Gamma_k>4$, the centered normalized alternating polynomial at $t_k$ is at least $1-2/\sqrt{\Gamma_k}>0$ because the two tails beginning at distance two alternate with decreasing magnitudes. In the finite positive overlaps, for $x=t_k$ or $v_k^d$, the terms at each fixed distance $j$ from the center converge after normalization to $e^{-j^2/(2c)}$. Indeed the terms at distance one have the required relative magnitudes by the definition of $t_k$ and (M2), and the consecutive logarithmic differences follow from $\log\Gamma_{k+j}$. Beyond any large fixed distance $L$, log-concavity bounds the outward ratios by a geometric factor $<1$ uniformly, using the ratio at $L$. Taking $L\to\infty$ therefore yields

\begin{equation*}
\frac{(-1)^kP(-x)}{a_kx^k}\longrightarrow
\sum_{j\in\mathbb Z}(-1)^j e^{-j^2/(2c)}
=\sqrt{2\pi c}\sum_{m\in\mathbb Z}e^{-2\pi^2 c(m+1/2)^2}>0.
\end{equation*}

Here the last identity is Poisson summation. \href{https://github.com/ProgramComputer/open-math-research/blob/f7e10f84dd21fec44ca58c60bae0879d12ebd89b/topics/balanced-arrays/orientation-normalization.md}{Orientation inequality}

If $k/d,(D-k)/d\to\infty$, set $w=v_k$. Then $n\mu/d^2,nh/d^2\to\infty$ from (M1). On the product-negative torus, character integration of the $n$th power of the row generating function gives $P(-w^d)$, as only the common column exponents survive. The two principal contributions have phase $\pm\Psi(w)$ and equal positive amplitude. Their sum differs from the exact polynomial by $o(1)$ times the amplitude of one contribution. This conclusion covers every tilt: use the small-mean proof for $\mu\le1/10$, the near-full proof for $h\le h_0$ with fixed absolute $h_0>0$, and the interior proof for $\mu\ge1/10$, $h\ge h_0$. All three estimates use the branch of the quadratic determinant obtained from the positive real tilt. At the first crossing $\cos\Psi(w)=(-1)^k$, giving the desired strict sign. \hyperref[sec:nearly-empty]{Nearly empty rows} \enspace \hyperref[sec:growing-cap-upper-saddle]{Near-full angular estimate} \enspace \hyperref[sec:interior-saddle]{Interior angular estimate}

For $k\le d$ or $D-k\le d$ use $t_k$, and at all other internal quotas use $v_k^d$. Inside their respective ranges the points increase by strict log-concavity and by continuity at the first crossings. At either join the adjacent geometric logarithms differ by $(\log\Gamma_k+\log\Gamma_{k+1})/2\to1$; replacing one geometric radius by the corresponding phase radius changes its logarithm by $o(1)$ by the overlap estimates. The endpoints are zero and positive infinity, with signs $1$ and $(-1)^D$. A subsequence with any failed interior sign has a further subsequence with limits of $k/d$ and $(D-k)/d$ in $[0,\infty]$. Their sum tends to infinity. The arguments above exclude a zero limit, a positive finite limit and two infinite limits. Thus the signs alternate at every one of these ordered points. By continuity the reflected polynomial has $D$ distinct positive zeros, exhausting its degree. This gives the stated simplicity and negativity and rules out the assumed sequence of counterexamples.
